\Askhsh{16805_SOLUTION}{1}
\begin{ylh}{1}
    \item [Ύλη από εικόνα]
\end{ylh}
\begin{wrapfigure}[9]{r}{5.5cm}
    \centering
    \begin{tikzpicture}[scale=0.3]
        % Ορισμός κορυφών
        \tkzDefPoint(0,12){D} % \Delta
        \tkzDefPoint(24,12){G} % \Gamma
        \tkzDefPoint(24,0){B}
        \tkzDefPoint(0,0){A}
        \tkzDefPoint(8,0){E}
        
        % Σχεδίαση του ορθογωνίου
        \tkzDrawPolygon(A,B,G,D)
        
        % Σχεδίαση των εσωτερικών τμημάτων
        \tkzDrawSegments(D,E G,E)
        
        % Σημεία (κόκκινες κουκκίδες)
        \tkzDrawPoints(A,B,G,D,E)
        
        % Ετικέτες
        \tkzLabelPoint[below left](A){A}
        \tkzLabelPoint[below right](B){B}
        \tkzLabelPoint[above right](G){$\Gamma$}
        \tkzLabelPoint[above left](D){$\Delta$}
        \tkzLabelPoint[below](E){E}
        
        % Ετικέτες μηκών
        \tkzLabelSegment[below](A,E){$x$}
        \tkzLabelSegment[below](E,B){$y$}
        \tkzLabelSegment[right](B,G){$z$}
    \end{tikzpicture}
\end{wrapfigure}
{\notoserif\bfseries\normalsize ΛΥΣΗ}
\begin{alist}
    \item Η περίμετρος του ορθογωνίου $\mathrm{AB}\Gamma\Delta$ είναι 72. Οπότε $2 \mathrm{AB} + 2 \mathrm{B}\Gamma = 72$ ή $2 (x + y) + 2 z = 72$ ή $x + y + z = 36$. Τα μήκη των τμημάτων $x, y, z$ είναι ανάλογα προς τους αριθμούς $2,4,3$ αντίστοιχα, άρα ισχύει
    \[ \frac{x}{2} = \frac{y}{4} = \frac{z}{3} = \frac{x+y+z}{2+4+3} = \frac{36}{9} = 4 \]
    Άρα $x = 2 \cdot 4 = 8$, $y = 4 \cdot 4 = 16$ και $z = 3 \cdot 4 = 12$.
    \item Από το Πυθαγόρειο θεώρημα στο τρίγωνο $\Gamma\mathrm{BE}$ έχουμε $\mathrm{\Gamma E}^2 = y^2 + z^2$ ή $\mathrm{\Gamma E}^2 = 16^2 + 12^2$, οπότε $\mathrm{\Gamma E}^2 = 400$ ή $\mathrm{\Gamma E} = 20$. Αντίστοιχα στο τρίγωνο $\Delta \mathrm{AE}$ για την υποτείνουσα $\Delta \mathrm{E}$ έχουμε $\Delta \mathrm{E}^2 = \mathrm{AE}^2 + \Delta \mathrm{A}^2$ ή $\Delta \mathrm{E}^2 = 8^2 + 12^2 = 208$, οπότε $\Delta \mathrm{E} = \sqrt{208} = \sqrt{16 \cdot 13} = 4\sqrt{13}$. Άρα η περίμετρος του τριγώνου $\Delta \mathrm{E}\Gamma$ ισούται με : $\Delta \mathrm{E} + \mathrm{E}\Gamma + \Delta \Gamma = 4\sqrt{13} + 20 + (8+16) = 44 + 4\sqrt{13}$.
\end{alist}